本文把四问统一为“有界观测形成可行域、覆盖圆给出安全清除点、有限空间覆盖证明无遗漏”。问题一给出了半平面交、状态判定、顶点枚举和直径计算算法，并以可实现的等边三角形反例证明直径圆不一定覆盖定位区域；清除时应改用最小包围圆。

问题二给出的候选域是三个 $1000\unit{m}$ 圆盘的交集，推荐点
$(761.429,\pm648.248)\unit{m}$ 同时满足保证接收和避免平行交会两项要求；参数扫描进一步显示，当 $R/r_0$ 逼近 2 时横向基线和最坏交会角会迅速退化。问题三采用 600 米方格顶点扫描、最近未完成点续扫和自适应定位，问题四通过凸包--半平面论证把发现保证扩展到未知方向的定向源。六轮官方演练 82 个源全部清除；同种子消融和定向比例压力分析补充验证了路线改进与保守定向策略。

当前策略优先保证完整性，移动距离仍是主要耗时，最近点续扫也不代表全局最短路线。正式测试前冻结核心代码；测试后仅据原始结果更新表格、摘要与结论，再复核日志和提交材料。正式结果若与合成趋势不一致，以正式原始日志为准。
